\section{Conclusion}

\label{sec:conclusion}
%% ============================================================================

Quasar achieves 204\,ms median finality over WAN (CPU BLS) and \textbf{10\,ms finality}
with GPU-accelerated BLS in co-located deployments, with simultaneous classical and
post-quantum cryptographic guarantees---a combination no other production protocol
provides. The dual-certificate mechanism adds less than 10\% overhead compared to
classical-only consensus, invalidating the assumption that post-quantum security
requires significant performance sacrifice.

At Lux's 1\,B gas block limit (\num{47619} txs/block), GPU BLS finality yields
\textbf{\num{4761904}\,TPS with finality}. CPU BLS provides a 100\,ms fallback
(\num{476190}\,TPS). At 100 validators over WAN, consensus-layer throughput reaches
\num{10400}\,TPS with O($n$) message complexity, positioning Quasar between
high-throughput DAG protocols (Narwhal+Tusk, $\sim$38{,}500\,TPS) and traditional
BFT protocols (Tendermint, $\sim$1{,}200\,TPS). For applications where finality
speed and quantum resistance matter more than raw throughput---cross-chain bridges,
financial settlement, regulatory-sensitive operations---Quasar is the optimal choice.

The benchmark data in this paper is reproducible. Hardware specifications, software
versions, and test configurations are fully documented. We encourage independent
reproduction and will publish the complete test harness as open source.

%% ============================================================================
%% REFERENCES
%% ============================================================================

\begin{thebibliography}{99}

\bibitem{quasar-paper}
Kelling, Z. (2025).
\textit{Quasar: Quantum-Secure Multi-Engine Consensus with Dual-Certificate Finality}.
Lux Industries Research.

\bibitem{hotstuff}
Yin, M., Malkhi, D., Reiter, M.K., Gueta, G.G., Abraham, I. (2019).
\textit{HotStuff: BFT Consensus with Linearity and Responsiveness}.
ACM PODC 2019. \texttt{arXiv:1803.05069}.

\bibitem{tendermint}
Buchman, E. (2016).
\textit{Tendermint: Byzantine Fault Tolerance in the Age of Blockchains}.
University of Guelph.

\bibitem{narwhal-tusk}
Danezis, G., Kokoris-Kogias, E., Sonnino, A., Spiegelman, A. (2022).
\textit{Narwhal and Tusk: A DAG-based Mempool and Efficient BFT Consensus}.
EuroSys 2022. \texttt{arXiv:2105.11827}.

\bibitem{casper-ffg}
Buterin, V., Griffith, V. (2017).
\textit{Casper the Friendly Finality Gadget}.
\texttt{arXiv:1710.09437}.

\bibitem{solana-wp}
Yakovenko, A. (2018).
\textit{Solana: A New Architecture for a High Performance Blockchain}.
Solana Labs Whitepaper v0.8.13.

\bibitem{snow-paper}
Team Rocket. (2020).
\textit{Scalable and Probabilistic Leaderless BFT Consensus through Metastability}.
\texttt{arXiv:1906.08936}.

\bibitem{ringtail}
NTT Research. (2024).
\textit{Pulsar: A Two-Round Post-Quantum Threshold Signature Scheme}.
Cryptology ePrint Archive.

\bibitem{fips204}
NIST. (2024).
\textit{FIPS 204: Module-Lattice-Based Digital Signature Standard (ML-DSA)}.
National Institute of Standards and Technology.

\bibitem{mysticeti}
Babel, K., Danezis, G., Kokoris-Kogias, E., Sonnino, A. (2024).
\textit{Mysticeti: Reaching the Latency Limits with Uncertified DAGs}.
Cornell University. \texttt{arXiv:2310.14821}.

\bibitem{blockbench}
Dinh, T.T.A., Wang, J., Chen, G., Liu, R., Ooi, B.C., Tan, K.L. (2017).
\textit{BLOCKBENCH: A Framework for Analyzing Private Blockchains}.
ACM SIGMOD 2017.

\bibitem{ethereum-pq}
Ethereum Foundation. (2025).
\textit{Post-Quantum Cryptography Research Roadmap}.
ethereum.org/research.

\bibitem{hotstuff2}
Malkhi, D., Nayak, K. (2023).
\textit{HotStuff-2: Optimal Two-Phase Responsive BFT}.
Cryptology ePrint Archive 2023/397.

\bibitem{bls-sigs}
Boneh, D., Lynn, B., Shacham, H. (2001).
\textit{Short Signatures from the Weil Pairing}.
ASIACRYPT 2001, LNCS 2248, pp. 514--532.

\bibitem{pbft}
Castro, M., Liskov, B. (1999).
\textit{Practical Byzantine Fault Tolerance}.
OSDI 1999.

\end{thebibliography}

%% ============================================================================
%% APPENDICES
%% ============================================================================

\appendix

